\documentclass{article}
\usepackage[T1]{fontenc}
\usepackage{lmodern}
\usepackage{graphicx}
\usepackage{amsmath,amssymb}
\usepackage[a4paper,margin=2cm]{geometry}
\usepackage[absolute,overlay]{textpos}
\setlength{\TPHorizModule}{1mm}
\setlength{\TPVertModule}{1mm}
\usepackage[indonesian]{babel}
\usepackage{hyperref}
\setlength{\emergencystretch}{2em}
% unit: Halaman 26

\begin{document}

% === Halaman 26 (python/native) ===

26

Cipher Block Chaining(CBC)

• Tujuan: membuat ketergantungan antar blok (mengatasi kelemahan mode ECB).

• Setiap blok cipherteks bergantung tidak hanya pada blok plainteksnya tetapi juga pada 

seluruh blok plainteks sebelumnya.  


• Hasil enkripsi blok sebelumnya di-umpan-balikkan ke dalam enkripsi blok yang current.

• Enkripsi blok pertama memerlukan blok semu (C0) yang disebut IV (initialization vector). 

• IV dapat diberikan oleh pengguna atau dibangkitkan secara acak oleh program.

• Pada dekripsi, blok plainteks diperoleh  dengan cara meng-XOR-kan IV dengan hasil 

dekripsi terhadap blok cipherteks pertama.

\end{document}
